\section{Conclusion}
\label{sec:conclusion}

We opened this paper with a concrete observation: controlling for MMLU general capability reduces the TruthfulQA $\times$ BBQ correlation from 0.732 to 0.343 --- more than half of what leaderboard comparisons might attribute to alignment coherence is, in fact, scale-driven confounding. Our work provides, to our knowledge, the first direct application of the Fisher $z$ raw-vs-partial difference test as a formal statistical diagnostic for this confound, and demonstrates it on $N{=}296$ open-weight LLMs with a pre-registered hypothesis and 57/57 test cases passing.

The three-stage verified mechanism is clear: (1) MMLU explains 49\% of TruthfulQA variance and 76\% of BBQ-proxy variance, confirming MMLU as a strong scale covariate; (2) controlling for MMLU reduces the pairwise alignment correlation by 53\% with Fisher $z$ $p = 3.22 \times 10^{-12}$, decisively confirming scale confounding; (3) the residual $\rho_{\text{partial}} = 0.343$ remains positive and significant ($p = 1.40 \times 10^{-9}$), indicating genuine alignment-specific co-movement beyond scale --- though the structural interpretation remains ambiguous at $N{=}296$.

\textbf{Future directions.} The most pressing next step is replication with genuine HELM Lite BBQ per-model accuracy scores. Extension to three-dimensional analysis (factuality $\times$ bias $\times$ safety) requires resolving HarmBench model name compatibility. Applying the same diagnostic to LLM LB v2 ($N > 1000$ models) would provide the statistical power to definitively resolve the scenario classification ambiguity.

\textbf{Methodological take-away.} The Fisher $z$ difference test between raw and partial Spearman is a practical, low-cost diagnostic that should become standard practice when interpreting alignment benchmark correlations. Capability control before interpreting benchmark co-movement is not optional --- it is methodologically necessary.
